\documentclass{article}
\usepackage{graphicx} % Required for inserting images

\usepackage{amsmath}
\usepackage{amsthm}
\usepackage{amssymb}
\usepackage{mathrsfs}
\usepackage{enumitem}
\newtheorem{theorem}{Theorem}
\newtheorem{lemma}{Lemma}
\newtheorem{corollary}{Corollary}
\theoremstyle{remark}\newtheorem{remark}{Remark}

% Theorem with manually specified number, to match Goldrei's numbering
\theoremstyle{plain}
\newtheorem{bookthminner}{Theorem}
\newenvironment{bookthm}[1]
  {\renewcommand\thebookthminner{#1}\bookthminner}
  {\endbookthminner}

\newcommand{\T}[0]{\mathcal{T}}
\newcommand{\B}[0]{\mathcal{B}}
\newcommand{\C}[0]{\mathcal{C}}
\newcommand{\R}[0]{\mathbb{R}}
\newcommand{\Q}[0]{\mathbb{Q}}
\newcommand{\Z}[0]{\mathbb{Z}}
\newcommand{\N}[0]{\mathbb{N}}
\newcommand{\e}[0]{\epsilon}

\begin{document}

2.38 Establish each of the following equivalences, where $\phi, \psi, \theta$, and all the $\theta_i$ are formulas.

(a) $((\phi \land \psi) \lor \neg\theta) \equiv ((\phi \lor \neg\theta) \land (\psi \lor \neg\theta))$

(b) $(\phi \to \neg\psi) \equiv (\psi \to \neg\phi)$

(c) $(\theta \to (\phi \lor \psi)) \equiv (\phi \lor (\psi \lor \neg\theta))$

(d) $((\theta_1 \land \theta_2) \land (\theta_3 \land \theta_4)) \equiv (\theta_1 \land ((\theta_2 \land \theta_3)) \land \theta_4))$

\begin{proof}[Proof of (a)]
    \begin{align*}
        ((\phi \land \psi) \lor \neg\theta)
            &\equiv (\neg\theta \lor (\phi \land \psi)) && \text{(by Theorem 2.3(b))}\\
        &\equiv \; ((\neg\theta \lor \phi) \land (\neg\theta \lor \psi)) && \text{(by Theorem 2.3(k) and transitivity)}
    \end{align*}
    
    But

    $(\neg\theta \lor \phi) \equiv (\phi \lor \neg\theta)$ and $(\neg\theta \lor \psi \equiv (\psi \lor \neg\theta)$

    by Theorem 2.3(b) so Exercise 2.37(b)

    $((\neg\theta \lor \phi) \land (\neg\theta \lor \psi)) \equiv ((\phi \lor \neg\theta) \land (\psi \lor \neg\theta))$.

    Then by transitivity of logical equivalence,

    $((\phi \land \psi) \lor \neg\theta) \equiv ((\phi \lor \neg\theta) \land (\psi \lor \neg\theta))$
    
    
\end{proof}

% ----------------------------

\begin{proof}[Proof of (b)]
    One of the logical equivalences Theorem 2.4(a) states is: $(\phi \to \psi) \equiv (\neg\psi \to \neg\phi)$.

    Treating the formulas above as metavariables and substituting $\phi \mapsto \phi, \psi \mapsto \neg\psi$ uniformly on both sides we arrive at
    $(\phi \to \neg\psi) \equiv (\neg\neg\psi \to \neg\phi)$.

    Theorem 2.3 (g) states the law of double negation which is $\neg\neg\phi \equiv \phi$. Substituting $\phi \mapsto \psi$ uniformly on both sides we arrive at $\neg\neg\psi \equiv \psi$.

    So then by exercise 2.37(d) $(\neg\neg\psi \to \neg\phi) \equiv (\psi \to \neg\phi)$ and transitively we get $(\phi \to \neg\psi) \equiv (\psi \to \neg\phi)$.
\end{proof}

% ----------------------------

\begin{proof}[Proof of (c)]
    One of the logical equivalences Theorem 2.4(a) states is: $(\phi \to \psi) \equiv (\neg\phi \lor \psi)$.

    Treating the above formulas as metavariables and substituting $\phi \mapsto \theta, \psi \mapsto (\phi \lor \psi)$ uniformly on both sides we arrive at $(\theta \to (\phi \lor \psi)) \equiv (\neg\theta \lor (\phi \lor \psi))$.

    By Theorem 2.3 (b) we have commutativity of $\lor$ so, $(\neg\theta \lor (\phi \lor \psi)) \equiv ((\phi \lor \psi) \lor \neg\theta)$.

    By Theorem 2.3 (f) we have associativity of $\lor$, so together with symmetry, $((\phi \lor \psi) \lor \neg\theta) \equiv (\phi \lor (\psi \lor \neg\theta))$.

    Lastly by transitivity we arrive at $(\theta \to (\phi \lor \psi)) \equiv (\phi \lor (\psi \lor \neg\theta))$.
\end{proof}

% ----------------------------

\begin{proof}[Proof of (d)]
    By Theorem 2.3 (e) we have $(\phi \land (\psi \land \theta)) \equiv ((\phi \land \psi) \land \theta)$.

    Then by symmetry we have $((\phi \land \psi) \land \theta) \equiv (\phi \land (\psi \land \theta))$.

    Treating the above formulas as metavariables and substituting $\phi \mapsto \theta_1, \psi \mapsto \theta_2, \theta \mapsto (\theta_3 \land \theta_4)$ uniformly on both sides we get $((\theta_1 \land \theta_2) \land (\theta_3 \land \theta_4)) \equiv (\theta_1 \land (\theta_2 \land (\theta_3 \land \theta_4)))$.

    By Theorem 2.3 (e) once again, substituting $\phi \mapsto \theta_2, \psi \mapsto \theta_3, \theta \mapsto \theta_4$ uniformly on both sides we get $(\theta_2 \land (\theta_3 \land \theta_4)) \equiv ((\theta_2 \land \theta_3) \land \theta_4)$. Then by Exercise 2.37(b) (with $\theta_1 \equiv \theta_1$, from Exercise 2.35(a)), $(\theta_1 \land (\theta_2 \land (\theta_3 \land \theta_4))) \equiv (\theta_1 \land ((\theta_2 \land \theta_3) \land \theta_4))$.

    Lastly by transitivity we arrive at $((\theta_1 \land \theta_2) \land (\theta_3 \land \theta_4)) \equiv (\theta_1 \land ((\theta_2 \land \theta_3) \land \theta_4))$.
\end{proof}

\end{document}

